\section{Introducción}

La Lengua de Señas Ecuatoriana (LSEC) constituye un componente fundamental del patrimonio lingüístico y cultural del país, siendo esencial para la comunicación de personas sordas y su integración en la sociedad. El estudio sistemático de la LSEC ha ganado importancia en años recientes, especialmente en el contexto de investigaciones en inteligencia artificial y procesamiento de lenguaje natural. Para realizar análisis rigurosos de videos de LSEC, la comunidad científica ha adoptado herramientas de anotación multimodal como ELAN (Eudico Language Annotator), desarrollada por el Max Planck Institute, que permite segmentar y etiquetar eventos lingüísticos en videos. Sin embargo, la arquitectura cerrada de ELAN, construida sobre la máquina virtual Java (JVM), presenta limitaciones significativas para la integración con modelos modernos de Inteligencia Artificial, que típicamente requieren entornos Python, librerías especializadas como PyTorch, y acceso a aceleración por GPU.

Este trabajo de integración curricular aborda la brecha de interoperabilidad entre ELAN y modelos de aprendizaje profundo mediante el desarrollo de un middleware de orquestación que actúa como intermediario. La solución adopta el patrón de diseño Sidecar, permitiendo que ELAN comunique solicitudes de procesamiento a un servicio independiente sin requerir modificaciones profundas en su arquitectura base. El middleware gestiona dinámicamente contenedores Docker para ejecutar modelos de segmentación y clasificación de videos, traduce protocolos heredados de ELAN (como AVATecH) a interfaces REST modernas, y orquesta la asignación de recursos computacionales (GPU/CPU) para viabilizar la ejecución de modelos complejos en computadoras estándar. Esta arquitectura garantiza la soberanía de los datos al mantener el procesamiento en el entorno local del usuario, sin dependencia de servicios en la nube, y preserva la estabilidad de ELAN mediante el aislamiento de componentes.

La implementación del middleware se estructuró en fases iterativas que avanzan desde un prototipo funcional con servicios base, pasando por gestión de modelos, runners de inferencia (dummy y nativo con PyTorch), procesamiento real de video con OpenCV, hasta integración final con ELAN mediante un bridge que traduce respuestas del middleware en anotaciones. El sistema propuesto resuelve la incompatibilidad Java-Python de forma elegante, permitiendo que investigadores trabajen con ELAN manteniendo acceso a modelos de IA de última generación para tareas como segmentación temporal de gestos, clasificación de glosas y extracción de keypoints corporales mediante visión computacional.

\subsection{Objetivo general}

Desarrollo de un middleware de orquestación de servicios de Inteligencia Artificial diseñado para integrar la herramienta de anotación multimodal ELAN con modelos modernos de aprendizaje profundo dedicados al procesamiento de videos de Lengua de Señas Ecuatoriana.

\subsection{Objetivos específicos}

Los objetivos específicos del proyecto se definen de la siguiente manera:

\begin{itemize}
\item Analizar el protocolo de extensión AVATecH de ELAN mediante técnicas de ingeniería inversa para definir los contratos de interfaz (CMDI) que permitan la inyección transparente de datos desde procesos externos hacia la línea de tiempo del lingüista, facilitando la comunicación bidireccional entre ELAN y el middleware.
\item Implementar el núcleo del orquestador de servicios utilizando Python con el framework FastAPI, desarrollando mecanismos de Inter-Process Communication (IPC) a través de APIs REST, gestor de colas de trabajos, y algoritmos de asignación dinámica de memoria para asegurar la estabilidad del sistema durante la inferencia de modelos pesados en GPU/CPU.
\item Desarrollar un pipeline de procesamiento que integre extracción de keypoints, segmentación temporal mediante modelos BiLSTM, clasificación de glosas con Transformer, y postprocesamiento de salidas de modelo para generar anotaciones semánticas compatibles con ELAN.
\item Validar el desempeño técnico del middleware mediante pruebas de carga en escenarios de anotación de LSEC, midiendo latencia de comunicación, consumo de recursos computacionales y la correcta recuperación de errores ante fallos de los modelos.
\end{itemize}

\subsection{Alcance}

El alcance del proyecto se centra en la definición, desarrollo y validación de la capa de infraestructura lógica del middleware de orquestación para ELAN. En la etapa inicial se realizó un análisis detallado de la especificación técnica de ELAN, su estructura de archivos y los protocolos nativos de extensión, con especial atención al protocolo AVATecH, para definir contratos de interfaz claros y reproducibles que permitan la inyección de datos desde procesos externos hacia la línea de tiempo del lingüista.

A continuación se implementó el middleware, enfocándose en mecanismos de comunicación interprocesos (IPC) mediante APIs REST, gestión de modelos con un registro centralizado, creación de runners especializados (dummy y nativo con PyTorch) y procesamiento de video real con bibliotecas de código abierto. La solución se diseñó de forma que la gestión de dependencias fuera agnóstica al sistema operativo del usuario final, garantizando portabilidad y estabilidad. Adicionalmente, se integraron componentes de postprocesamiento temporal y traducción de respuestas para generar anotaciones compatibles con ELAN sin producir archivos EAF automáticos desde el middleware.

Finalmente, el sistema se validó en un entorno controlado con pruebas de caja blanca y caja negra. Se comprobó la orquestación de modelos sobre videos de LSEC, se midieron métricas de desempeño como uso de CPU/RAM, latencia y tiempos de respuesta, y se evaluó la tolerancia a fallos bajo condiciones de degradación. Se documentó la arquitectura final para asegurar la mantenibilidad y escalabilidad del componente. El proyecto no incluye el desarrollo de modelos de IA desde cero, integración con servicios en la nube, generación automática de archivos EAF, expansión a herramientas de anotación distintas a ELAN, ni entrenamiento de modelos. El desarrollo se limita a un entorno local, priorizando interoperabilidad Java-Python, gestión dinámica de recursos de hardware y viabilidad en computadoras estándar con o sin GPU.



-------------------- poner en metodología
Para el desarrollo de esta solución se aplicó un enfoque ágil e iterativo, donde se optó por hacer uso de Scrum como marco de trabajo ágil para organizar ciclos de revisión y mejora continua. En este contexto, la metodología utilizada combinó la naturaleza iterativa del desarrollo incremental con principios ágiles de adaptación permanente lo que permitió realizar ajustes frecuentes en función de los resultados de análisis y pruebas realizadas a lo largo del desarrollo.